\begin{pblock}{Reproducibility}
Every product is a cloud-optimised GeoTIFF with a STAC item beside it carrying
the run's resolved parameters, so provenance travels with the file rather than
living in a notebook. The toolkit is deliberately granular: no
\texttt{run\_everything()}, one function per scientific step.
\end{pblock}

\vspace{4mm}
{\color{PSlate}\small
Bishop-Taylor, R.\ et al.\ (2019). Between the tides: modelling the elevation
of Australia's exposed intertidal zone at continental scale. \emph{Est.\
Coast.\ Shelf Sci.} 223, 115--128.\\[1.2mm]
Hart-Davis, M.\ et al.\ (2021). EOT20: a global ocean tide model from
multi-mission satellite altimetry. \emph{ESSD} 13, 3869--3884.\\[1.2mm]
Schramm, M.\ et al.\ (2021). The openEO API. \emph{ISPRS Int.\ J.\
Geo-Inf.} 10(3), 125.\\[1.2mm]
Granadeiro, J.\,P.\ et al.\ (2021). Using Sentinel-2 images to estimate
topography, tidal-stage lags and exposure periods over large intertidal
areas. \emph{Remote Sens.} 13(2), 320.

\vspace{3mm}
Contains modified Copernicus Sentinel-2 data (2016--2025) via the Copernicus
Data Space openEO backend. Tides: EOT20 (SEANOE), GOT4.10 (NASA GSFC) via
pyTMD. Field validation: RTK GNSS campaign, Ría de Villaviciosa.
\par}
